\documentclass[12pt, a4paper]{article}
\usepackage[UTF8]{ctex}
\usepackage{amsmath, amssymb}
\usepackage{geometry}
\geometry{margin=1in}
\pagestyle{empty}
\begin{document}
\begin{center}
    {\LARGE\bfseries 线性代数 第 3 次作业}\\[0.5em]
    {\large 线性代数（A）}\\[0.3em]
    李生丹 \quad | \quad 2026 年 10 月
\end{center}
\hrule
\vspace{1em}

\noindent\textbf{题 1.} 计算下列矩阵的行列式：
\[A = \begin{pmatrix} 2 & 1 & 3 \\ 4 & 5 & 6 \\ 7 & 8 & 9 \end{pmatrix}\]

\vspace{0.8em}
\noindent\textbf{题 2.} 求下列矩阵的逆矩阵：
\[A = \begin{pmatrix} 2 & 1 \\ 5 & 3 \end{pmatrix}\]

\vspace{0.8em}
\noindent\textbf{题 3.} 求矩阵的特征值与特征向量：
\[A = \begin{pmatrix} 4 & 1 \\ 2 & 3 \end{pmatrix}\]

\vspace{0.8em}
\noindent\textbf{题 4.} 解线性方程组：
\[\left[\begin{array}{ccc|c} 1 & 2 & 3 & 6 \\ 2 & 4 & 6 & 12 \\ 1 & 1 & 1 & 3 \end{array}\right]\]

\vspace{0.8em}
\noindent\textbf{题 5.} 求下列矩阵的秩：
\[A = \begin{pmatrix} 1 & 2 & 3 \\ 2 & 4 & 6 \\ 3 & 6 & 9 \end{pmatrix}\]

\vspace{0.8em}
\noindent\textbf{题 6.} 已知矩阵
\[A = \begin{pmatrix} 1 & 2 \\ 3 & 4 \end{pmatrix} , \quad B = \begin{pmatrix} 2 & 0 \\ 1 & 3 \end{pmatrix}\]
计算矩阵乘积 $AB$。

\vspace{0.8em}
\noindent\textbf{题 7.} 用 Gram-Schmidt 方法将下列向量组正交化：
\[\alpha_1 = \begin{pmatrix} 1 \\ 1 \\ 0 \end{pmatrix} , \quad \alpha_2 = \begin{pmatrix} 1 \\ 0 \\ 1 \end{pmatrix}\]

\vspace{0.8em}
\noindent\textbf{题 8.} 判断矩阵 $A$ 是否可对角化；若可以，求出对角化所需的矩阵 $P$ 与对角矩阵 $D$：
\[A = \begin{pmatrix} 4 & 1 \\ 2 & 3 \end{pmatrix}\]

\end{document}
